\begin{choice}[topic={充分不必要条件的应用}]
若 ``$x=2$'' 是 ``$m^2x^2 - (m+3)x + 4 = 0$'' 的充分不必要条件，则实数 $m$ 的值为
\begin{tasks}(2)
	\task 1
	\task $-\frac{1}{2}$
	\task $-\frac{1}{2}$ 或 1
	\task -1 或 $-\frac{1}{2}$
\end{tasks}
\end{choice}
\begin{choice-sol}
\textbf{答案 B}

命题``$p$ 是 $q$ 的充分不必要条件''包含两层含义，必须同时满足：
\begin{itemize}
	\item \textbf{充分性}: $p \Rightarrow q$。
	\item \textbf{不必要性}: $q \not\Rightarrow p$。
\end{itemize}
设 $p: x=2$，$q: m^2x^2 - (m+3)x + 4 = 0$。

\textbf{检验充分性}:
将 $x=2$ 代入方程 $q$ 中，方程必须成立.
\[ 4m^2 - 2(m+3) + 4 = 0 \implies 2m^2 - m - 1 = 0 \]
因式分解得 $(2m+1)(m-1) = 0$，
解得 $m=1$ 或 $m=-\frac{1}{2}$。

\textbf{检验不必要性}:
方程 $q$ 的解集 $S$ 必须真包含 $\{2\}$，即 $S \neq \{2\}$。
\begin{itemize}
	\item 当 $m=1$ 时，方程为 $x^2 - 4x + 4 = 0$，
	解集 $S=\{2\}$.
	此时 $q \Rightarrow p$ 成立，不满足``不必要''条件，故舍去.
	\item 当 $m=-\frac{1}{2}$ 时，
	方程为 $\frac{1}{4}x^2 - \frac{5}{2}x + 4 = 0$，
	即 $x^2 - 10x + 16 = 0$.
	解集 $S=\{2, 8\}$.
	此时 $S \neq \{2\}$，满足``不必要''条件。
\end{itemize}
综上，唯一满足条件的实数 $m$ 的值为 $-\frac{1}{2}$。
\end{choice-sol}
